\newpage
\section{Probability Theory and Stochastic Processes}
In this chapter, we introduce the basic notation and fundamental tools of probability theory and measure theory that will be used throughout the subsequent chapters. The section is organized as follows: 
\begin{enumerate}
    \item probability space and random variables;
    \item modes of convergence;
    \item the Law of Large Numbers (LLN) and the Central Limit Theorem (CLT);
    \item martingales, Markov chains and Brownian motion. 
\end{enumerate}

This section mainly follows Introduction to Probability Theory (Tsinghua University), Stochastic Processes (Peking University), Stochastic Analysis (Peking University), Math C218A (UC Berkeley), Math C218B (UC Berkeley) and the book Probability: Theory and Examples by Rick Durrett~\cite{durrett2019probability}, Sinho Chewi's notes of MATH C218A, and Introduction to Stochastic Processes~\cite{lawler2018introduction}.

\subsection{Probability Space and Random Variables}
\subsubsection{Probability space}
A probability space \( (\Omega, \mathcal{F}, P) \) contains three elements:
\begin{itemize}
    \item The space \(\Omega\): this is a non-empty set. It can be viewed as the set of all possible outcomes.
    \item The \(\sigma\)-field \(\mathcal{F}\): this can be viewed as a collection of all the events.
    \item The probability measure \(P\): this is a function from \(\mathcal{F}\) to \([0, 1]\). It gives a probability to each event.
\end{itemize}

\begin{defn}[$\sigma-$Field]
    Suppose \(\mathcal{F}\) is a non-empty collection of subsets of \(\Omega\).
    \begin{itemize}
        \item It is a field if it is closed under \textbf{complementation} and closed under \textbf{union}:
        \[
        A \in \mathcal{F} \implies A^c \in \mathcal{F}, \quad A_1, A_2 \in \mathcal{F} \implies A_1 \cup A_2 \in \mathcal{F}.
        \]
        \item It is a \textbf{monotone class} if
        \[
        A_j \in \mathcal{F}, \, A_j \subset A_{j+1}, \, 1 \leq j < \infty \implies \bigcup_j A_j \in \mathcal{F},
        \]
        and
        \[
        A_j \in \mathcal{F}, \, A_j \supset A_{j+1}, \, 1 \leq j < \infty \implies \bigcap_j A_j \in \mathcal{F}.
        \]
        \item It is a \(\sigma\)-field if it is closed under \textbf{complementation} and closed under \textbf{countable union}:
        \[
        A \in \mathcal{F} \implies A^c \in \mathcal{F}, \quad A_j \in \mathcal{F}, \, 1 \leq j < \infty \implies \bigcup_j A_j \in \mathcal{F}.
        \]
    \end{itemize}
\end{defn}

\begin{defn}[Generated $\sigma$-Fields]
Given any collection \( C \) of sets, the \(\sigma\)-field generated by \( C \) is the intersection of all \(\sigma\)-fields containing \( C \).
\end{defn}

\begin{defn}[Probability Measure]
    Suppose \(\mathcal{F}\) is a \(\sigma\)-field on \(\Omega\). A probability measure \(P\) is a function from \(\mathcal{F}\) to \([0, 1]\) satisfying the following axioms:
    \begin{itemize}
        \item \(P[E] \geq 0\) for all \(E \in \mathcal{F}\);
        \item \(P[\Omega] = 1\);
        \item If \(\{E_j\}_{j}\) is a countable collection of pairwise disjoint sets in \(\mathcal{F}\), then:
        \[
        P\left[\bigcup_j E_j\right] = \sum_j P[E_j].
        \]
    \end{itemize}
    These axioms imply the following consequences:
    \begin{itemize}
        \item \(P[E^c] = 1 - P[E]\);
        \item \(P[E \cup F] + P[E \cap F] = P[E] + P[F]\);
        \item Continuity: if \(E_n \uparrow E\) or \(E_n \downarrow E\), then \(P[E_n] \to P[E]\);
        \item \(P\left[\bigcup_j E_j\right] \leq \sum_j P[E_j]\).
    \end{itemize}
\end{defn}

\begin{rem}
    We care about the $\sigma$-field because that is where our random variables live. It defines the rules under which stochastic analysis can work. Sometimes we need to enlarge our space to a bigger $\sigma$-field and make sure the measure still fits — this is guaranteed by the following theorem.
\end{rem}


\begin{thm}[Carathéodory’s Extension Theorem]
    Suppose \(\mathcal{F}_0\) is a field and \(\mathcal{F}\) is the \(\sigma\)-field generated by \(\mathcal{F}_0\). Suppose \(\mu\) is a probability measure on \(\mathcal{F}_0\). Then there exists a unique probability measure on \(\mathcal{F}\) that coincides with \(\mu\) on \(\mathcal{F}_0\).
\end{thm}

\begin{rem}
    The construction: 
    \[
    \tau(A)=\inf\left\{
    \sum_{n=1}^\infty \mu(B_n): B_n\in \mathcal{F}, \ A\subset\bigcup_{n=1}^\infty B_n
    \right\}.
    \]
\end{rem}



\subsubsection{Random Variables}
\begin{defn}[Random Variables]
    A real-valued random variable is a function \( X: \Omega \to \mathbb{R} \) such that
    \[
    X^{-1}(B) \in \mathcal{F}, \quad \forall B \in \mathcal{B}.
    \]
    In other words, a random variable is just a \textbf{measurable} function from \( (\Omega, \mathcal{F}) \) to \( (\mathbb{R}, \mathcal{B}) \).
\end{defn}
\begin{defn}[Distribution]
    Each random variable \( X \) induces a probability measure \( \mu \) on \( (\mathbb{R}, \mathcal{B}) \) by the following correspondence:
    \[
    \mu[B] = P[X^{-1}(B)] = P[X \in B], \quad \forall B \in \mathcal{B}.
    \]
    The measure \( \mu \) is called the law (or the distribution) of \( X \), denoted by \( \mathcal{L}(X) \); its associated distribution function is called the distribution function of \( X \), denoted by \( F_X \).
\end{defn}
The concept of ``expectation'' is the same as integration in the probability space \( (\Omega, \mathcal{F}, \mathbb{P}) \).
\begin{cor}
    If \( X \) takes only positive integer values, we have
    \[
    \mathbb{E}[X] = \sum_{n=1}^{\infty} \mathbb{P}[X \geq n].
    \]
\end{cor}

\begin{defn}[$p$-th Moment]
    For any \( p \in (0, \infty) \), define
    \[
    L^p(\Omega, \mathcal{F}, \mathbb{P}) = \left\{ X \text{ random variable on } (\Omega, \mathcal{F}, \mathbb{P}) : \mathbb{E}[|X|^p] < \infty \right\}.
    \]
    For \( X \in L^p \), we call \( \mathbb{E}[|X|^p] \) the \( p \)-th moment of \( X \).
\end{defn}

\begin{defn}[Independent]
    The random variables \( \{ X_j, 1 \leq j \leq n \} \) are independent if, for any Borel sets \( \{ B_j, 1 \leq j \leq n \} \), we have
    \[
    \mathbb{P}\left[\bigcap_{j=1}^n \{ X_j \in B_j \}\right] = \prod_{j=1}^n \mathbb{P}[X_j \in B_j].
    \]
    The random variables \( \{ X_j, j \geq 1 \} \) are independent if \( \{ X_j, 1 \leq j \leq n \} \) are independent for all \( n \).
\end{defn}

\begin{thm}
    Suppose \( \{ A_j, 1 \leq j \leq n \} \) are independent and each \( A_j \) is a \(\pi\)-system. Denote by \( \sigma(A_j) \) the \(\sigma\)-field generated by \( A_j \) for each \( j \). Then \( \{ \sigma(A_j), 1 \leq j \leq n \} \) are independent.
\end{thm}

\begin{exam}[Kolmogorov’s 0-1 Law]
    Let \( \{ X_n \} \) be a sequence of independent random variables. Let
    \[
    G_n = \sigma(X_k, k \geq n) \quad \text{and} \quad G_\infty = \bigcap_{n \geq 1} G_n.
    \]
    Then \( G_\infty \) is trivial, i.e., for any \( A \in G_\infty \), we have
    \[
    \mathbb{P}[A] = 0 \text{ or } 1.
    \]
    Define \( S_n = \sum_{j=1}^{n} X_j \). It is clear that the following events are in \( G_\infty \):
    \[
    \lim_{n \to \infty} S_n \text{ exists}, \quad \lim_{n \to \infty} \frac{S_n}{n} \text{ exists}, \quad \limsup_{n \to \infty} \frac{S_n}{n} > 0.
    \]
    Whereas, the following event is not in \( G_\infty \):
    \[
    \limsup_{n \to \infty} S_n > 0.
    \]
\end{exam}

\begin{proof}
    On the one hand, we have \( A \in G_{n+1} \) for any \( n \). Thus, \( A \) is independent of \( \sigma(X_1, \dots, X_n) \) for any \( n \). Therefore, \( A \) is independent of \( \sigma(X_n, n \geq 1) \). On the other hand, \( A \) is measurable with respect to \( \sigma(X_n, n \geq 1) \). Therefore, \( A \) is independent of itself. This implies \( \mathbb{P}[A] = \mathbb{P}[A \cap A] = \mathbb{P}[A]^2 \), thus \( \mathbb{P}[A] \in \{ 0, 1 \} \).
\end{proof}

Now let's talk about Gaussians. The key thing to know is that Gaussians stay Gaussian under linear combinations and projections.
\begin{prop}
   The following are equivalent:
    \begin{enumerate}
        \item The random vector \( \xi = (\xi_1, \dots, \xi_d)^\top \) follows a \( d \)-dimensional Gaussian distribution.
        \item For any \( a_1, \dots, a_d \), the linear combination \( \sum_{k=1}^d a_k \xi_k \) follows a one-dimensional Gaussian distribution.
    \end{enumerate} 
\end{prop}


\subsection{Convergence}
\subsubsection{Convergence: Almost Sure, in Probability, and in \texorpdfstring{$L^p$}{Lp}}
\begin{defn}[Almost sure convergence (a.s.))]
    The sequence of random variables \( \{ X_n \} \) converges a.s. to the random variable \( X \) if there exists a null set \( \mathcal{N} \) such that
    \[
    \lim_{n \to \infty} X_n(\omega) = X(\omega), \quad \forall \omega \in \Omega \setminus \mathcal{N}.
    \]
\end{defn}

\begin{defn}[Convergence in probability] 
    The sequence \( \{ X_n \} \) converges in probability to the random variable \( X \) if, for every \( \epsilon > 0 \), we have
    \[
    \lim_{n \to \infty} \mathbb{P}[|X_n - X| > \epsilon] = 0.
    \]
\end{defn}

\begin{defn}[Convergence in \( L^p \))] 
    Assume \( p \geq 1 \). The sequence \( \{ X_n \} \) converges in \( L^p \) to the random variable \( X \) if \( X_n \in L^p \), \( X \in L^p \) and
    \[
    \lim_{n \to \infty} \mathbb{E}[|X_n - X|^p] = 0.
    \]
\end{defn}

It is straightforward to prove that almost sure convergence (\(a.s.\)) implies convergence in probability, and convergence in \(L^p\) implies convergence in probability. However, almost sure convergence (\(a.s.\)) and \(L^p\) may not imply each other. Now we talk about an example: using probability methods to prove Weierstrass Theorem.

\begin{exam}[Weierstrass Theorem]
    Let \( f \) be a continuous function on \( [0, 1] \). Define the polynomial:
\[
f_n(x) = \sum_{j=0}^{n} \binom{n}{j} x^j (1 - x)^{n - j} f\left( \frac{j}{n} \right).
\]
This is called the Bernstein polynomial of degree \( n \) associated to \( f \). Then we have
\[
\sup_{x \in [0, 1]} |f_n(x) - f(x)| \to 0, \quad n \to \infty.
\]
\end{exam}

\begin{proof}
    Suppose \( \{ X_j, j \geq 1 \} \) are i.i.d. Bernoulli variables with parameter \( p \in [0, 1] \): \( \mathbb{P}[X_j = 1] = p, \mathbb{P}[X_j = 0] = 1 - p \), and set \( S_n = \sum_{j=1}^{n} X_j \). Then we find
\[
\mathbb{E}[X_j] = p, \quad \text{var}(X_j) = p(1 - p), \quad \mathbb{E}[S_n] = np, \quad \text{var}(S_n) = np(1 - p).
\]
Moreover, 
\[
\mathbb{P}[S_n = j] = \binom{n}{j} p^j (1 - p)^{n - j}, \quad \text{thus} \quad f_n(p) = \mathbb{E}[f(S_n/n)].
\]
Note that \( f \) is bounded: \( |f| \leq M \), and it is uniformly continuous: for any \( \epsilon > 0 \), there exists \( \delta > 0 \) such that \( |f(x) - f(y)| \leq \epsilon \) as long as \( |x - y| \leq \delta \). Thus,
\[
|f_n(p) - f(p)| \leq \mathbb{E}[|f(S_n/n) - f(p)|] \leq \epsilon + 2M \mathbb{P}[|S_n/n - p| > \delta].
\]
Note that
\[
\mathbb{P}[|S_n/n - p| > \delta] \leq \frac{\text{var}(S_n)}{n^2 \delta^2} = \frac{p(1 - p)}{n \delta^2} \leq \frac{1}{4n \delta^2}.
\]
Thus,
\[
|f_n(p) - f(p)| \leq \epsilon + \frac{M}{2n \delta^2}.
\]
\end{proof}

We give a useful lemma to prove almost surely convergence.
\begin{thm}[Borel-Cantelli Lemma]
\label{BC_lemma}
\begin{itemize}
    \item For arbitrary sequence \( \{ E_n \} \), we have
    \[
    \sum_n \mathbb{P}[E_n] < \infty \quad \Rightarrow \quad \mathbb{P}[E_n \text{ i.o.}] = 0.
    \]
    \item If the events \( \{ E_n \} \) are independent, we have
    \[
    \sum_n \mathbb{P}[E_n] = \infty \quad \Rightarrow \quad \mathbb{P}[E_n \text{ i.o.}] = 1.
    \]
\end{itemize}
\end{thm}

\subsubsection{Weak Convergence}
\begin{defn}[Weak Convergence]
    A sequence of measures \( \{ \mu_n \} \) converges weakly to a measure \( \mu \) if
    \[
    \mu_n((a, b]) \to \mu((a, b]), \quad \text{for all continuity points} \ a, b \ \text{of} \ \mu.
    \]
    We denote by \( \mu_n \Rightarrow \mu \).
\end{defn}

\begin{prop}
    Suppose \( \{ \mu_n \} \) and \( \mu \) are probability measures. Then \( \mu_n \Rightarrow \mu \) if and only if
\[
\lim_{n \to \infty} \int f(x) \, \mu_n[d x] \to \int f(x) \, \mu[d x], \quad \forall f \in C_b:=\text{(bounded continuous functions)}.
\]
\end{prop}

\subsubsection{Convergence in Distribution}
\begin{defn}[Convergence in Distribution]
    A sequence of random variables \( \{ X_n \} \) converges in distribution to a random variable \( X \) if distribution function \( \mathcal{L}(X_n) \Rightarrow \mathcal{L}(X) \). We denote by \( X_n \xrightarrow{d} X \), or \( X_n \to X \) in distribution.
\end{defn}
Convergence in probability implies convergence in distribution. Furthermore, convergence in distribution is equivalent to the convergence of \textbf{characteristic functions}.

\begin{defn}[Characteristic Function]
For any random variable $X$ with distribution $\mu$, its characteristic function is defined to be:
\[
f: \mathbb{R} \to \mathbb{C}, \quad f(t) := \mathbb{E}\left[e^{itX}\right] = \int e^{itx} \, \mu[dx].
\]
\end{defn}

The distribution is uniquely identified by the characteristic function.
\begin{thm}
Suppose $f$ is the characteristic function for the probability measure $\mu$. For $x < y$, we have
\[
\mu[(x, y)] + \frac{1}{2}\mu[\{x\}] + \frac{1}{2}\mu[\{y\}] = \lim_{T \to \infty} \frac{1}{2\pi} \int_{-T}^{T} \frac{e^{-itx} - e^{-ity}}{it} f(t) \, dt.
\]
\end{thm}

\begin{thm}
Let $\{\mu_n\}$ be a sequence of probability measures with characteristic functions $\{f_n\}$. Suppose that
\begin{enumerate}
    \item $f_n$ converges everywhere in $\mathbb{R}$ and defines the limiting function $f$;
    \item $f$ is continuous at $t = 0$.
\end{enumerate}
Then we have
\begin{enumerate}
    \item $\mu_n \xrightarrow{d} \mu$ where $\mu$ is a probability measure;
    \item the characteristic function of $\mu$ is $f$.
\end{enumerate}
\end{thm}

At the end of this subsection, we introduce some useful convergence results. Let \( (X_n) \) be a sequence of random variables and let \( X \) be a random variable such that \( X_n \to X \) (a.s.):
\[
\mathbb{P}(X_n \to X) = 1.
\]

We state the convergence properties as follows:

\begin{itemize}
    \item \textbf{(MON)} If \( 0 \leq X_n \uparrow X \), then \( \mathbb{E}(X_n) \leq \mathbb{E}(X) < \infty \);
    \item \textbf{(FATOU)} If \( X_n \geq 0 \), then \( \mathbb{E}(X_n) \leq \liminf \mathbb{E}(X_n) \);
    \item \textbf{(DOM)} If \( |X_n(\omega)| \leq Y(\omega), \, \forall (n, \omega), \text{ and } \mathbb{E}(Y) < \infty \), then
    \[
    \mathbb{E}(|X_n - X|) \to 0,
    \]
    \item \textbf{(SCHEFFÉ)} If \( \mathbb{E}(|X_n|) \to \mathbb{E}(|X|) \), then
    \[
    \mathbb{E}(|X_n - X|) \to 0;
    \]
    \item \textbf{(BDD)} If there exists a constant \( K \), such that \( |X_n(\omega)| \leq K, \, \forall (n, \omega) \), then
    \[
    \mathbb{E}(|X_n - X|) \to 0.
    \]
\end{itemize}


\subsection{The Law of Large Numbers and the Central Limit Theorem}
\subsubsection{Strong Law of Large Numbers}
We will give three versions (4th moment SLLN, 2nd moment SLLN and SLLN).
\begin{thm}[4th Moment SLLN]
Let $(X_i, 1 \leq i < \infty)$ be IID, $EX = 0$, and $EX^4 < \infty$. Write $S_n = \sum_{i=1}^{n} X_i$. Then
\begin{enumerate}
    \item $ES_n^4 \leq 3n^2 EX^4$
    \item $S_n/n \to 0$ as $n \to \infty$.
\end{enumerate}
\end{thm}

If $EX = \mu$, applying the theorem to $X - \mu$ shows that $S_n/n \to \mu$ a.s.

\begin{proof}
    \begin{align*}
    ES_n^4 &=\sum_i \sum_j \sum_k \sum_l E[X_i X_j X_k X_l]\\
    &= nEX^4 + \binom{4}{2}\binom{n}{2} E[X_1^2 X_2^2] \\
    &= nEX^4 + 3n(n-1) \underbrace{(EX^2)^2}_{\leq EX^4}
    \end{align*}
    since $(EY)^2 \leq E(Y^2)$.
 Fix $\varepsilon > 0$.
    \begin{align*}
    P\left(\left|\frac{S_n}{n}\right| \geq \varepsilon\right) &\leq E\left|\frac{S_n}{n}\right|^4 \cdot \frac{1}{\varepsilon^4} \\
    &\leq \varepsilon^{-4} n^{-4} \cdot 3n^2 EX^4 \\
    &\leq 3\varepsilon^{-4} EX^4 n^{-2}
    \end{align*}
    This implies that
    \[
    \sum_n P\left(\left|\frac{S_n}{n}\right| \geq \varepsilon\right) \leq \sum_n 3\varepsilon^{-4} EX^4 n^{-2} < \infty
    \]
    By Borel-Cantelli Lemma~\ref{BC_lemma}, $S_n/n \to 0$ a.s. We used the fact that $s^4 = |s|^4$ and $s^2 = |s|^2$, but this does not work for the third moment: $s^3 \neq |s|^3$.
\end{proof}

\begin{thm}[2nd Moment SLLN]
Given $(X_i, 1 \leq i < \infty)$, with $EX_i \equiv 0$, let $\sup_i EX_i^2 = B < \infty$ and the $X_i$ be \textbf{orthogonal}, $E(X_i X_j) = 0$, $j \neq i$. (We are not assuming independence!) Write
\[
S_n = \sum_{i=1}^{n} X_i
\]
Then $S_n/n \to 0$ a.s.
\end{thm}

\begin{proof}
Since $\operatorname{var}(S_n) \leq nB$, Chebyshev's inequality implies
\[
P\left(\frac{|S_n|}{n} \geq \varepsilon\right) \leq \frac{nB}{n^2 \varepsilon^2} = \frac{B}{n\varepsilon^2}
\]
Take $n(j) = j^2$.
\[
P\left(\left|\frac{S_{n(j)}}{n(j)}\right| \geq \varepsilon\right) \leq \frac{B}{\varepsilon^2} \frac{1}{j^2}
\]
Use Borel-Cantelli~\ref{BC_lemma}.
\[
\frac{S_{n(j)}}{n(j)} \to 0 \quad \text{a.s.} \quad \text{as } j \to \infty
\]
It is enough to prove $D_j / j^2 \to 0$ a.s., for
\[
D_j = \max_{j^2 \leq n < (j+1)^2} |S_n - S_{j^2}|
\]
Then
\[
D_j^2 = \max_{j^2 \leq n < (j+1)^2} (S_n - S_{j^2})^2
\]
\[
ED_j^2 \underset{\text{crude}}{\leq} \sum_{n=j^2}^{(j+1)^2 - 1} E(S_n - S_{j^2})^2
\]
Since
\[
E(S_n - S_{j^2})^2 = \operatorname{var}\left(\sum_{j^2+1}^{n} X_i\right) \leq B(n - j^2)
\]
Letting $n = j^2 + i$, we have
\[
ED_j^2 \leq B \sum_{i=1}^{2j+1} i = \frac{1}{2}(2j+1)(2j+2)B
\]
We have
\[
P\left(\frac{D_j}{j^2} \geq \varepsilon\right) \leq \frac{ED_j^2}{\varepsilon^2 j^4} \in O(j^{-2})
\]
Borel-Cantelli Lemma~\ref{BC_lemma} implies that $D_j / j^2 \to 0$ as $j \to \infty$.
\end{proof}

\begin{thm}[SLLN]
\label{SLLN}
Let $(X_i)$ be IID with $E|X| < \infty$. Then $S_n/n \to EX$ a.s. as $n \to \infty$.
\end{thm}

First prove some lemmas. We'll use~\ref{SLLN_tech2} and truncate + center techniques.
\begin{thm}
\label{SLLN_tech1}
Let $(X_i)$ be independent, with $EX_i = 0$ and $\sigma_i^2 = \operatorname{var}(X_i) < \infty$. If $\sum_{i=1}^{\infty} \sigma_i^2 < \infty$, then $\sum_{i=1}^{\infty} X_i$ converges a.s.
\end{thm}

\begin{proof}
Define $M_k = \sup_{n > k} \left|\sum_{i=k+1}^{n} X_i\right|$. It is enough to show that $M_k \to 0$ a.s. as $k \to \infty$. Define also $W_k = \sup_{n_2 > n_1 > k} \left|\sum_{i=n_1+1}^{n_2} X_i\right|$ and note that $M_k \leq W_k \leq 2M_k$ and $W_k$ decreases as $k$ increases.
\[
P\left(\sup_{k < n \leq N} \left|\sum_{i=k+1}^{N} X_i\right| \geq \varepsilon\right) \underset{martingale \ maximal \ ineq.}{\leq} \varepsilon^{-2} \operatorname{var}\left(\sum_{i=k+1}^{N} X_i\right) = \varepsilon^{-2} \sum_{i=k+1}^{N} \sigma_i^2
\]

Taking $N \to \infty$, $P(M_k > \varepsilon) \leq \varepsilon^{-2} \sum_{i=k+1}^{\infty} \sigma_i^2$.
\[
P(W_k > \varepsilon) \leq P\left(M_k > \frac{\varepsilon}{2}\right) \leq 4\varepsilon^{-2} \sum_{i=k+1}^{\infty} \sigma_i^2 \to 0 \text{ as } k \to \infty
\]

Taking $k \to \infty$, then $W_k \downarrow W_\infty$ for some $W_\infty$ a.s. Then $P(W_\infty > \varepsilon) = 0$, which implies that $W_\infty = 0$ a.s., which implies that $W_k \downarrow 0$ a.s. and $M_k \to 0$ a.s.
\end{proof}

\begin{lem}[Deterministic Lemma (Kronecker)]
\label{Kronecker}
Let $(x_n)$ be a sequence of reals, $S_n = \sum_{i=1}^{n} x_i$, $0 < a_n \uparrow \infty$ as $n \uparrow \infty$. If $\sum_i x_i / a_i$ converges, then $S_n / a_n \to 0$.
\end{lem}

\begin{cor}
\label{SLLN_tech2}
Let $(X_i)$ be independent, $EX_i = 0$, $EX_i^2 < \infty$, and $S_n = \sum_{i=1}^{n} X_i$. If $0 < a_n \uparrow \infty$ as $n \uparrow \infty$ and if $\sum_n EX_n^2 / a_n^2 < \infty$, then $S_n / a_n \to 0$ a.s.
\end{cor}

\begin{proof}
Theorem~\ref{SLLN_tech1} implies that $\sum_n X_n / a_n$ converges a.s. Then Lemma~\ref{Kronecker} implies that $S_n / a_n \to 0$ a.s.
\end{proof}

\begin{proof}[Proof of Theorem~\ref{SLLN}]
The idea is to truncate, center, and then apply 9.4.

If $Z \geq 0$, then
\[
EZ^k = \int_0^{\infty} kz^{k-1} P(Z \geq z) \, dz.
\]

Define $Y_k = X_k 1_{(|X_k| \leq k)}$. Then
\[
\sum_k P(Y_k \neq X_k) = \sum_{k=1}^{\infty} P(|X| > k) \leq \int_0^{\infty} P(|X| > x) \, dx = E|X| < \infty.
\]

Then Borel-Cantelli Lemma~\ref{BC_lemma} implies that $P(Y_k = X_k, \text{ ultimately}) = 1$. It is enough to prove that $(1/n) \sum_{k=1}^{n} Y_k \to EX$ a.s.

Center: define $X_k' = Y_k - EY_k$. \textit{Claim:} $\sum_k \operatorname{var}(X_k')/k^2 < \infty$.
\begin{align*}
EY_k^2 &= \int_0^{\infty} 2y P(|Y_k| > y) \, dy = \int_0^{\infty} \underbrace{2y P(k \geq |X_k| \geq y) 1_{(y \leq k)}}_{\text{Check this!}} \, dy \\
&\leq \int_0^{\infty} 2y P(|X_k| \geq y) 1_{(y \leq k)} \, dy
\end{align*}
\begin{align*}
\sum_k \frac{\operatorname{var}(X_k)}{k^2} &\leq \sum_k \frac{EY_k^2}{k^2} \leq \sum_k \frac{1}{k^2} \int_0^{\infty} 2y P(|X| \geq y) 1_{(y \leq k)} \, dy \\
&= \int_0^{\infty} \underbrace{\left(\sum_k \frac{1}{k^2} 1_{(y \leq k)} 2y\right)}_{G(y)} P(|X| \geq y) \, dy
\end{align*}

\textit{Claim:} $G(y) \leq 4$, for all $0 < y < \infty$. Since $G(y) \leq \sum_k 1/k^2 \leq 2$ for $y \leq 1$, this is true for $y \leq 1$. Take $y > 1$.
\[
\frac{1}{k^2} \leq \int_{k-1}^{k} \frac{1}{x^2} \, dx
\]
so
\[
\sum_k \frac{1}{k^2} 1_{(y \leq k)} = \sum_{k \geq \lceil y \rceil} \frac{1}{k^2} \leq \int_{\lceil y \rceil - 1}^{\infty} \frac{1}{x^2} \, dx = \frac{1}{\lceil y \rceil - 1}
\]

Since $y > 1$,
\[
G(y) \leq \frac{2y}{\lceil y \rceil - 1} \leq 4.
\]
Then
\[
\sum_k \frac{\operatorname{var}(X_k')}{k^2} \leq 4 \int_0^{\infty} P(|X| \geq y) \, dy = 4E|X| < \infty
\]

Apply~\ref{SLLN_tech2} to $(X_n')$: $(1/n) \sum_{i=1}^{n} X_i' \to 0$ a.s., so $(1/n) \sum_{i=1}^{n} (Y_i - EY_i) \to 0$ a.s. Note that
\[
EY_i = EX 1_{(|X| \leq i)} \to EX
\]
as $i \to \infty$. By dominated convergence, $(1/n) \sum_{i=1}^{n} (EY_i - EX) \to 0$ a.s. Add the two equations to get $(1/n) \sum_{i=1}^{n} (Y_i - EX) \to 0$ a.s., which implies that $(1/n) \sum_{i=1}^{n} Y_i \to EX$ a.s.
\end{proof}

\subsubsection{Central Limit Theorem}
\begin{thm}[IID Central Limit Theorem]
Let $(X_i, i \geq 1)$ be IID, $EX = \mu$, $\operatorname{var}(X) = \sigma^2 < \infty$. Then,
\[
\frac{S_n - n\mu}{\sqrt{n}} \xrightarrow{d} \operatorname{Normal}(0, \sigma^2).
\]
\end{thm}

\begin{proof}
WLOG take $\mu = 0$. It is enough to show
\[
\underbrace{\phi_{S_n/\sqrt{n}}(t)}_{\text{Left}} \to \exp\left(-\frac{\sigma^2 t^2}{2}\right).
\]

Also,
\[
\phi_{S_n/\sqrt{n}}(t) = \left(\phi_X\left(\frac{t}{\sqrt{n}}\right)\right)^n = \left(1 + \frac{n(\phi_X(t/\sqrt{n}) - 1)}{n}\right)^n.
\]

It is enough to show $n(\phi_X(t/\sqrt{n}) - 1) \to \sigma^2 t^2 / 2$. The bound for $n = 2$ and $EX = 0$ is
\[
\left|\phi_X(s) - \left(1 - \frac{s^2 \sigma^2}{2}\right)\right| = o(s^2).
\]

Then, with $s = t/\sqrt{n}$,
\[
\text{Left} = n\left(\frac{t^2}{n}\frac{\sigma^2}{2} + o\left(\frac{t^2}{n}\right)\right) = \frac{t^2 \sigma^2}{2} + n \cdot o\left(\frac{t^2}{n}\right) \to \frac{t^2 \sigma^2}{2}. \qedhere
\]
\end{proof}

\subsubsection{Law of the Iterated Logarithm}
\begin{thm}[Law of the Iterated Logarithm]
Let $\{X_n\}$ be independent, identically distributed random variables with zero means and unit variances. Let $S_n = X_1 + \cdots + X_n$. Then
\[
\limsup_{n \to \infty} \frac{|S_n|}{\sqrt{2n \log \log n}} = 1 \quad \text{a.s.}
\]
\end{thm}

\subsection{Stochastic Processes}
\begin{defn}[Usual Hypothesis]
$(\Omega, \mathcal{F}, (\mathcal{F}_t)_{t \geq 0}, P)$ is a probability space equipped with a filtration, satisfying the usual conditions:
\begin{enumerate}
    \item $\mathcal{F}_{t_1} \subset \mathcal{F}_{t_2}$ for $t_1 \leq t_2$, and $\mathcal{F}_t \subset \mathcal{F}$;
    \item $(\Omega, \mathcal{F}, P)$ is a complete probability space;
    \item $\mathcal{F}_0$ contains all $P$-null sets in $\mathcal{F}$;
    \item (Right-continuity) For all $t > 0$, $\displaystyle \mathcal{F}_t = \bigcap_{s > t} \mathcal{F}_s$.
\end{enumerate}
\end{defn}

\begin{defn}[Stopping Time]
A random variable $T : \Omega \longrightarrow [0, \infty]$ is a \textit{stopping time} of the filtration $(\mathcal{F}_t)$ if $\{T \leq t\} \in \mathcal{F}_t$, for every $t \geq 0$. The \textit{$\sigma$-field of the past before $T$} is then defined by
\[
\mathcal{F}_T = \{A \in \mathcal{F}_\infty : \forall t \geq 0, \, A \cap \{T \leq t\} \in \mathcal{F}_t\}.
\]
\end{defn}

\begin{rem}
    Easy to verify that $\mathcal{F}_T$ is indeed a $\sigma$-field. We have $\EE [X_\tau] = \EE X_0$ for stopping time $\tau$ and martingale $X_{\cdot}$ under some assumptions. 
\end{rem}

\begin{defn}[Adaptedness]
A stochastic process $(X_t)_{t \in T}$ is said to be \textit{adapted} to the filtration $(\mathcal{F}_t)_{t \in T}$ if, for every $t \in T$, $X_t$ is a random variable on $(\Omega, \mathcal{F}_t)$, i.e., for all $a \in \mathbb{R}$,
\[
\{\omega : X_t(\omega) \leq a\} \in \mathcal{F}_t.
\]
\end{defn}

\subsubsection{Markov Chains}
\begin{defn}[Markov Process]
Consider a filtered probability space $(\Omega, \mathcal{F}, (\mathcal{F}_t)_{t \in T}, P)$, where $(X_t)_{t \in T}$ is an adapted stochastic process with respect to $(\mathcal{F}_t)_{t \in T}$. We say $(X_t)_{t \in T}$ is a \textit{Markov process} with respect to $(\mathcal{F}_t)_{t \in T}$ if for every bounded Borel function $f$ and every $t > s$,
\[
E[f(X_t) \mid \mathcal{F}_s] = E[f(X_t) \mid \sigma(X_s)] \Big(= E[f(X_t) \mid X_s]\Big).
\]
\end{defn}

\begin{rem}
    Define $p(s,x;t,y)$ as the transition probability function. Easy to get Chapman-Kolmogorov Equation:
    \[
    \int_{\mathbb{R}^n} f(y) P(s, x; t, \mathrm{d}y) = \int_{\mathbb{R}^n} \left( \int_{\mathbb{R}^n} f(y) P(r, z; t, \mathrm{d}y) \right) P(s, x; r, \mathrm{d}z).
    \]
\end{rem}

\begin{thm}[The MC Convergence Theorem]
Suppose the chain is \textbf{irreducible} and \textbf{positive-recurrent}, so the stationary $\pi$ exists. If the chain is also aperiodic, then $\mathbb{P}_{\mu_0}(X_n = j) \xrightarrow{n \to \infty} \pi(j) \; \forall j \; \forall \mu_0$.
\end{thm}
Using SLLN, we can easily get the following results.
\begin{thm}[Ergodic Theorem for Markov Chains]
Consider an irreducible, positive-recurrent MC. Let $\pi$ be the stationary distribution. Take $f : S \to \mathbb{R}$ such that $\sum_x \pi(x)|f(x)| < \infty$. Then,
\[
\frac{1}{t} \sum_{n=1}^{t} f(X_n) \xrightarrow{a.s.} \bar{f} := \sum_x \pi(x) f(x) \qquad \text{as } t \to \infty.
\]
\end{thm}

\begin{rem}
    Besides, some interesting techniques in the analysis of Markov processes are deeply connected to combinatorics. Using
    \[
    \sum_{n\ge 0} \frac{\binom{2n}{n}}{2^{2n}} \asymp \sum_{n\ge 0}\frac{1}{\sqrt{n}} \to \infty, \quad \sum_{n\ge 0} \frac{\binom{2n}{n}^2}{4^{2n}} \asymp \sum_{n\ge 0}\frac{1}{n} \to \infty,
    \]
    we can get 1D, 2D random walks are recurrent. Another example is more complicated.
\end{rem}

\begin{exam}[Catalan Number]
    Consider a simple symmetric random walk $(S_n)$ starting at $S_0=0$.
Let $\tau_i = \inf\{n \geq 1: S_n = i\}$.
We derive the number of paths of length $2n$ from $0$ to $0$ that stay $\geq 0$.

For $n+i$ even, the set $\{S_n = i\}\setminus\{\tau_i = n\} = \{\tau_i < n,\, S_n = i\}$ consists of paths that visited $i$ before time $n$.
Splitting by $S_{n-1}$ and applying reflection at time $\tau_i$, the contributions from $S_{n-1}=i+1$ and $S_{n-1}=i-1$ are equal, giving
\[
  P(\tau_i < n,\, S_n = i) = P(S_{n-1} = i+1) = \frac{n-i}{n}\,P(S_n = i).
\]
Therefore
\[
  P(\tau_i = n) = \frac{i}{n}\,P(S_n = i).
\]
Set $i=1$ and $n=2k+1$.  The event $\{\tau_1 = 2k+1\}$ forces $S_1,\ldots,S_{2k}\leq 0$ with $S_{2k}=0$ (the last step is necessarily $+1$). The number of such paths of the first $2k$ steps is
\[
  \frac{1}{2k+1}\binom{2k+1}{k+1}.
\]
Reflecting $S\mapsto -S$ counts instead the paths from $0$ to $0$ of length $2k$ staying $\geq 0$.  Simplifying:
\[
  \frac{1}{2k+1}\binom{2k+1}{k+1}
  = \frac{(2k)!}{(k+1)!\,k!}
  = \frac{1}{k+1}\binom{2k}{k}
  = C_k.
\]
\end{exam}

Another interesting topic is \textbf{Continuous-Time Markov Chains}. The continuous-time viewpoint is important in many problems, as it brings in the tools of ordinary differential equations.
\begin{exam}[Poisson Processes]
    A counting process $(X_t)_{t\ge0}$ with $X_0=0$ is a \emph{Poisson process of rate $\lambda$} if it has independent, stationary increments and
\[
  \mathbb P\{X_{t+\Delta t} - X_t = 1\} = \lambda\Delta t + o(\Delta t), \qquad
  \mathbb P\{X_{t+\Delta t} - X_t \ge 2\} = o(\Delta t).
\]
 
\paragraph{Approach 1 (binomial limit)}
Write $X_t = \sum_{j=1}^{n}(X_{jt/n}-X_{(j-1)t/n})$.  For large $n$ the summands are independent $\mathrm{Bernoulli}(\lambda t/n)$, so
\[
  \mathbb P\{X_t=k\} = \lim_{n\to\infty}\binom{n}{k}\Bigl(\frac{\lambda t}{n}\Bigr)^{\!k}\Bigl(1-\frac{\lambda t}{n}\Bigr)^{\!n-k} = e^{-\lambda t}\frac{(\lambda t)^k}{k!}.
\]
 
\paragraph{Approach 2 (ODE)}
Set $P_k(t)=\mathbb P\{X_t=k\}$.  The defining conditions give 
\[
P_k'(t)=\lambda P_{k-1}(t)-\lambda P_k(t).
\]
Substituting $f_k(t)=e^{\lambda t}P_k(t)$ reduces this to $f_k'=\lambda f_{k-1}$, $f_k(0)=0$.  Since $f_0\equiv1$, induction yields $f_k(t)=(\lambda t)^k/k!$.

\paragraph{Approach 3 (interarrival times)}
Let $T_n$ be the time between the $(n{-}1)$th and $n$th arrival.  The $T_i$ are i.i.d.\ and memoryless: $\mathbb P\{T_i\ge s+t\mid T_i\ge s\}=\mathbb P\{T_i\ge t\}$, so $T_i\sim\mathrm{Exp}(\lambda)$.  (Match rates: $\mathbb E[Y_n]=n/\lambda$ should give $\approx\!\lambda t$ arrivals by time $t$; or directly, $\mathbb P\{T_1>t\}=\mathbb P\{X_t=0\}=e^{-\lambda t}$.)  Then $Y_n=T_1+\cdots+T_n\sim\mathrm{Gamma}(n,\lambda)$, and
\[
  \mathbb P\{X_t=k\}=\mathbb P\{Y_k\le t < Y_{k+1}\}=\int_0^t \frac{\lambda^k s^{k-1}}{(k{-}1)!}e^{-\lambda s}\cdot\lambda e^{-\lambda(t-s)}\,ds = e^{-\lambda t}\frac{(\lambda t)^k}{k!}.
\]
\end{exam}

\subsubsection{Martingale}
\begin{defn}[Martingale]
Let $(X_t)_{t \in T}$ be an adapted stochastic process on $(\Omega, \mathcal{F}, (\mathcal{F}_t)_{t \in T}, P)$. If for all $s < t$,
\[
E[X_t \mid \mathcal{F}_s] = X_s,
\]
then $(X_t)_{t \in T}$ is called a \textit{martingale} with respect to $(\mathcal{F}_t)_{t \in T}$, or $(X_t, \mathcal{F}_t)_{t \in T}$ is called a martingale. When the filtration $(\mathcal{F}_t)_{t \in T}$ is clear from context, we simply call $(X_t)_{t \in T}$ a martingale.

If $E[X_t \mid \mathcal{F}_s] \geq X_s$, then $(X_t, \mathcal{F}_t)_{t \in T}$ is called a \textit{submartingale}; if $E[X_t \mid \mathcal{F}_s] \leq X_s$, then $(X_t, \mathcal{F}_t)_{t \in T}$ is called a \textit{supermartingale}.
\end{defn}

\begin{defn}[Doob-Meyer Decomposition, rough version]
Let $(X_t)_{t \geq 0}$ be an $(\mathcal{F}_t)_{t \geq 0}$-submartingale with RCLL paths. If $\{X_t\}_{t \geq 0}$ satisfies certain regularity conditions, then $(X_t)_{t \geq 0}$ can be uniquely decomposed as
\[
X_t = M_t + C_t, \quad t \geq 0
\]
where $(M_t)_{t \geq 0}$ is an $(\mathcal{F}_t)_{t \geq 0}$-martingale, and $(C_t)_{t \geq 0}$ is a predictable, right-continuous, increasing process with $E\, C_t < \infty$, $\forall t \geq 0$, and $C_0 = 0$. $(C_t)_{t \geq 0}$ is called the \textit{compensator} of $(X_t)_{t \geq 0}$.
\end{defn}

\begin{rem}
    Easy to verify that if $M_t$ is a martingale, then $M_t^2$ is a submartingale. Use $\langle M\rangle_t$ to denote the $C_t$ and define the stochastic integral on general continuous martingale.
\end{rem}

A useful result (which we also use for the proof of SLLN) is as follows.
\begin{thm}[Doob's Martingale Maximal Inequality]
Let $(\xi_t, \mathcal{F}_t)_{t \geq 0}$ be a martingale with RCLL paths (right-continuous with left limits). If for some $p \geq 1$ and all $t \geq 0$, $E|\xi_t|^p < \infty$, then for any $T < \infty$ and $a > 0$,
\[
P\left(\sup_{t \in [0,T]} |\xi_t| \geq a\right) \leq \frac{E|\xi_T|^p}{a^p}.
\]
For any $a, b > 0$,
\[
P\left(\sup_{t \in [0,T]} \xi_t \geq a\right) \leq \frac{E\, e^{b\xi_T}}{e^{ba}}.
\]
\end{thm}

\begin{thm}[Doob's Submartingale Convergence Theorem]
Let $(\xi_n, \mathcal{F}_n)_{n \geq 0}$ be a submartingale with $\sup_n E|\xi_n| < \infty$. Then $\lim_{n \to \infty} \xi_n = \xi_\infty$ almost everywhere, and $E|\xi_\infty| < \infty$.
\end{thm}

\subsubsection{Itô Calculus}
\begin{defn}[Brownian Motion]
Let $(B_t)_{t \geq 0}$ be a stochastic process. If it satisfies the following three conditions:
\begin{enumerate}
    \item $B_{t+s} - B_t$ follows a normal distribution $N(0, \sigma^2 s)$;
    \item For any $0 < t_1 < t_2 < \cdots < t_n$, $B_0, B_{t_1} - B_0, \cdots, B_{t_n} - B_{t_{n-1}}$ are mutually independent;
    \item For all $\omega$, $B_t(\omega)$ is continuous in $t$.
\end{enumerate}
Then $(B_t)_{t \geq 0}$ is called a (one-dimensional) \textit{Brownian motion}. In these notes, we always assume that the Brownian motion satisfies $\sigma^2 = 1$; when $B_0 = 0$, it is called a \textit{standard Brownian motion}.
\end{defn}
\begin{rem}[History of Brownian Motion]
    \textbf{Einstein's derivation (1905).} Consider $n$ particles in a fluid. Assume: 
    \begin{enumerate}
        \item particles move independently;
        \item increments over non-overlapping time intervals are independent.
    \end{enumerate}
    Let $\varphi(\Delta)$ be the symmetric density of displacement over a time interval $\tau$, with $\varphi(\Delta) = \varphi(-\Delta)$.

Let $f(x,t)$ be the particle density. Then
\[
f(x, t+\tau) = \int_{-\infty}^{+\infty} f(x-\Delta, t)\,\varphi(\Delta)\,\mathrm{d}\Delta.
\]
Expand both sides: the left in $\tau$, the right in $\Delta$. By symmetry of $\varphi$, odd-order terms vanish. Keeping only leading terms and setting $D = \frac{1}{\tau}\int \frac{\Delta^2}{2}\,\varphi(\Delta)\,\mathrm{d}\Delta$, we obtain
\[
\frac{\partial f}{\partial t} = D\frac{\partial^2 f}{\partial x^2}.
\]
With initial condition $f(x,0) = n\,\delta(x)$, the solution is
\[
f(x,t) = \frac{n}{\sqrt{4\pi Dt}}\,e^{-\frac{x^2}{4Dt}}.
\]
The ratio $f(x,t)/n$ is precisely the density of a Brownian motion. In particular, the root-mean-square displacement scales as $\sqrt{2Dt} \sim \sqrt{t}$.
\end{rem}

\begin{prop}[Kolmogorov backward/Forward equation]
Let $\{p(t, x, y), \, x \in \mathbb{R}^n, \, y \in \mathbb{R}^n\}$ be the transition density of an $n$-dimensional standard Brownian motion. Then
\begin{equation}
\label{Kol_eq}
\frac{\partial p(t, x, y)}{\partial t} = \frac{1}{2} \Delta_x p(t, x, y) = \frac{1}{2} \sum_{i=1}^{n} \frac{\partial^2 p(t, x, y)}{\partial x_i^2}
\end{equation}
\begin{equation}
\label{FP_eq}
\frac{\partial p(t, x, y)}{\partial t} = \frac{1}{2} \Delta_y p(t, x, y) = \frac{1}{2} \sum_{i=1}^{n} \frac{\partial^2 p(t, x, y)}{\partial y_i^2}
\end{equation}
Equation~\ref{Kol_eq} is called the \textit{Kolmogorov backward equation} (the first equation), and~\ref{FP_eq} is called the \textit{Kolmogorov forward equation} (the second equation), also known as the \textit{Fokker--Planck equation}.
\end{prop}

\begin{prop}[Martingale Properties]
For a standard $(\mathcal{F}_t)_{t \geq 0}$-Brownian motion $(B_t)_{t \geq 0}$,
\begin{enumerate}
    \item $(B_t)_{t \geq 0}$ is an $(\mathcal{F}_t)_{t \geq 0}$-martingale;
    \item $(B_t^2 - t)_{t \geq 0}$ is an $(\mathcal{F}_t)_{t \geq 0}$-martingale;
    \item $(z_t = e^{cB_t - \frac{c^2 t}{2}})_{t \geq 0}$ is an $(\mathcal{F}_t)_{t \geq 0}$-martingale.
\end{enumerate}
\end{prop}

\begin{rem}
    Based on the martingale properties, we define stochastic integral using $\dd B_t$, which is well-known Itô Integral.
\end{rem}

\begin{defn}[Itô Integral]
For a stochastic process $(\phi_t)_{t \geq 0}$ in $\mathscr{L}_T^2$, let $(\phi_t^{(n)})_{t \geq 0}$ be a sequence of predictable simple (step) processes approximating it with respect to $(\mathcal{F}_t)_{t \geq 0}$, satisfying
\[
\|\phi^{(n)} - \phi\|^2 = E \int_0^T \left|\phi_t^{(n)} - \phi_t\right|^2 \mathrm{d}t \to 0, \quad \text{as } n \to \infty.
\]
(At this point $\int_0^T \phi_t^{(n)} \, \mathrm{d}B_t$ is a Cauchy sequence in $L^2(\Omega)$.) Denote the limit of $\int_0^T \phi_t^{(n)} \, \mathrm{d}B_t$ in $L^2(\Omega)$ by $\int_0^T \phi_t \, \mathrm{d}B_t$, i.e.,
\[
\int_0^T \phi_t \, \mathrm{d}B_t := L^2(\Omega)\text{-}\lim_{n \to \infty} \int_0^T \phi_t^{(n)} \, \mathrm{d}B_t,
\]
which is called the \textit{It\^{o} stochastic integral} of $(\phi_t)_{0 \leq t \leq T}$ with respect to $(B_t)_{t \geq 0}$.
\end{defn}
\begin{rem}
    Different choices of representative points will influence the result. We use left endpoints in the definition of the Itô integral. If we use midpoints instead, we obtain the Stratonovich integral.
\end{rem}
\begin{prop}[Properties of Stochastic Integrals]
.
\begin{enumerate}
    \item If $\tau$ is an $(\mathcal{F}_t)_{t \geq 0}$-stopping time with $\tau \leq T$, then
    \[
    \int_0^{\tau} \phi_t \, \mathrm{d}B_t = \int_0^{T} \phi_t \mathbf{1}_{(0,\tau]}(t) \, \mathrm{d}B_t.
    \]

    \item $\left\{\xi_t \overset{\mathrm{def}}{=} \int_0^t \phi_s \, \mathrm{d}B_s\right\}$ is an $(\mathcal{F}_t)_{t \in [0,T]}$-martingale.

    \item $\left\{\eta_t \overset{\mathrm{def}}{=} \left(\int_0^t \phi_s \, \mathrm{d}B_s\right)^2 - \int_0^t \phi_s^2 \, \mathrm{d}s\right\}$ is an $(\mathcal{F}_t)_{t \in [0,T]}$-martingale, and
    \[
    E\left[\left(\int_s^t \phi_u \, \mathrm{d}B_u\right)^2 \bigg| \mathcal{F}_s\right] = E\left[\int_s^t \phi_u^2 \, \mathrm{d}u \bigg| \mathcal{F}_s\right].
    \]

    \item If $(\phi_t)_{t \in [0,T]}$ is a bounded adapted process with respect to $(\mathcal{F}_t)_{t \in [0,T]}$, i.e., $\forall t$, $|\phi_t| < M < \infty$, then
    \[
    \left\{\zeta_t \overset{\mathrm{def}}{=} e^{\int_0^t \phi_s \, \mathrm{d}B_s - \frac{1}{2}\int_0^t \phi_s^2 \, \mathrm{d}s}\right\}
    \]
    is an $(\mathcal{F}_t)_{t \in [0,T]}$-martingale.
\end{enumerate}
\end{prop}

\begin{thm}[Itô's Formula]
Let $F(t, x)$ be once continuously differentiable in $t$ and twice continuously differentiable in $x$. Then
\begin{align*}
F(t, B_t) - F(0, B_0) &= \int_0^t \frac{\partial F(s, B_s)}{\partial s} \, \mathrm{d}s + \int_0^t \frac{\partial F(s, B_s)}{\partial x} \, \mathrm{d}B_s + \frac{1}{2} \int_0^t \frac{\partial^2 F(s, B_s)}{\partial x^2} \, \mathrm{d}s \\
&= \int_0^t F_t'(s, B_s) \, \mathrm{d}s + \int_0^t F_x'(s, B_s) \, \mathrm{d}B_s + \frac{1}{2} \int_0^t F_{xx}''(s, B_s) \, \mathrm{d}s.
\end{align*}
In differential form, this is written as
\[
\mathrm{d}F(t, B_t) = F_t'(t, B_t) \, \mathrm{d}t + \frac{1}{2} F_{xx}''(t, B_t) \, \mathrm{d}t + F_x'(t, B_t) \, \mathrm{d}B_t.
\]
\end{thm}
\begin{rem}
    A useful trick to remember:
    \[
    (\dd B_t)^2 = \dd t.
    \]
\end{rem}

\subsubsection{Stochastic Differential Equations}

At last, we talk about Stochastic Differential Equations (SDE). Consider a general form:
\[
\mathrm{d}\xi_t = \mathbf{b}(t, \xi_t) \, \mathrm{d}t + \Sigma(t, \xi_t) \, \mathrm{d}\mathbf{B}_t.
\]

\begin{thm}[Feynman-Kac Formula]
Let
\begin{equation}
\mathrm{d}\xi_t = \mathbf{b}(\xi_t) \, \mathrm{d}t + \Sigma(\xi_t) \, \mathrm{d}\mathbf{B}_t
\end{equation}
and denote $\mathcal{L} = \frac{1}{2}\sum_{i,j=1}^{n} a_{ij} \frac{\partial^2}{\partial x_i \partial x_j} + \sum_{i=1}^{n} b_i(x) \frac{\partial}{\partial x_i}$, where $(a_{ij})_{i,j=1,\cdots,n} = \Sigma\Sigma^\top$. Let $f \in C_0^2(\mathbb{R}^n)$, $q(x) \in C(\mathbb{R}^n)$, and $q(x)$ be bounded below.
\begin{enumerate}
    \item Let $v(t, x) = E\left(f(\xi_t) \exp\left(-\int_0^t q(\xi_s) \, \mathrm{d}s\right) \bigg| \xi_0 = x\right)$. Then
    \begin{equation}
    \label{eq:FK_formula}
    \begin{cases}
    \dfrac{\partial v}{\partial t} = \mathcal{L}v - qv & t > 0, \; x \in \mathbb{R}^n \\[6pt]
    v(0, x) = f(x) & x \in \mathbb{R}^n
    \end{cases}.
    \end{equation}
    \item If $w(t, x) \in C^{1,2}(\mathbb{R} \times \mathbb{R}^n)$ is bounded on $K \times \mathbb{R}^n$, $K \subset \mathbb{R}$, and $w$ is a solution to~\eqref{eq:FK_formula}, then $w(t, x) = v(t, x)$.
\end{enumerate}
\end{thm}

\begin{exam}[Diffusion Model]
    Consider the forward OU process
\[
\mathrm{d}x_t = -\tfrac{1}{2}x_t\,\mathrm{d}t + \mathrm{d}B_t, \quad x_0 \sim p_0.
\]
A direct computation gives $x_t | x_0 \sim N(\alpha_t x_0,\, \sigma_t^2 I)$ where $\alpha_t = e^{-t/2}$ and $\sigma_t^2 = 1 - e^{-t}$. As $t \to \infty$, the distribution $p_t$ converges exponentially to $N(0, I)$ --- structure is destroyed and replaced by pure noise.

By the Fokker--Planck equation, the density $p(x,t)$ satisfies
\[
\frac{\partial p}{\partial t} = \nabla \cdot \left(\tfrac{1}{2}x\, p\right) + \tfrac{1}{2}\Delta p = \tfrac{1}{2}\nabla \cdot \left(x\, p + \nabla p\right).
\]

\paragraph{Reverse SDE.} Anderson (1982) showed that for a general forward SDE $\mathrm{d}x_t = f(x,t)\,\mathrm{d}t + g(t)\,\mathrm{d}B_t$, a reverse process is given by
\[
\mathrm{d}\tilde{x}_t = \left[f(\tilde{x}_t, t) - g^2(t)\,\nabla_x \log p(\tilde{x}_t, t)\right]\mathrm{d}t + g(t)\,\mathrm{d}\bar{B}_t, \quad t: T \to 0,
\]
where $\bar{B}_t$ is a backward Brownian motion.
\begin{proof}
The forward SDE $\mathrm{d}x_t = f\,\mathrm{d}t + g\,\mathrm{d}B_t$ has Fokker--Planck equation
\[
\frac{\partial p}{\partial t} = -\nabla \cdot (f\, p) + \tfrac{1}{2}g^2 \Delta p.
\]
Set $s = T - t$ and $q(x,s) = p(x, T-s)$. Since $\partial q / \partial s = -\partial p / \partial t$, the density $q$ must satisfy
\[
\frac{\partial q}{\partial s} = \nabla \cdot (f\,p) - \tfrac{1}{2}g^2 \Delta p. \tag{$\star$}
\]
Now we verify that the proposed reverse SDE, with drift $h = -f + g^2 \nabla_x \log p$ and diffusion $g$, reproduces $(\star)$. Its Fokker--Planck equation in $s$ reads $\partial q / \partial s = -\nabla \cdot (h\,q) + \tfrac{1}{2}g^2 \Delta q$. Substituting $q = p$ and using $p\,\nabla \log p = \nabla p$:
\begin{align*}
-\nabla \cdot (h\,p) + \tfrac{1}{2}g^2 \Delta p
&= \nabla \cdot (f\,p) - \underbrace{\nabla \cdot (g^2 \nabla p)}_{=\, g^2 \Delta p} + \tfrac{1}{2}g^2 \Delta p
= \nabla \cdot (f\,p) - \tfrac{1}{2}g^2 \Delta p,
\end{align*}
which matches $(\star)$.
\end{proof}

For our OU process, $f(x,t) = -x/2$ and $g = 1$, so the reverse SDE reads
\[
\mathrm{d}\tilde{x}_t = \left[-\tfrac{1}{2}\tilde{x}_t - \nabla_x \log p(\tilde{x}_t, t)\right]\mathrm{d}t + \mathrm{d}\bar{B}_t.
\]
If we knew the score function $\nabla_x \log p(x,t)$, we could simulate this backward from $\tilde{x}_T \sim N(0,I)$ and recover samples from $p_0$. The entire problem reduces to \textbf{learning the score}.

\paragraph{Denoising score matching} We want to minimize $L_t(\theta) = \mathbb{E}_{x \sim p_t}\left[\|s_\theta(x,t) - \nabla_x \log p(x,t)\|^2\right]$, but $\nabla_x \log p(x,t)$ is unknown. The key trick is integration by parts. Expand the square and note that the cross term satisfies
\[
\int \langle s_\theta,\, \nabla_x p\rangle\,\mathrm{d}x = -\int (\nabla \cdot s_\theta)\, p\,\mathrm{d}x.
\]
Now use $p(x,t) = \int p_t(x|x_0)\,p_0(x_0)\,\mathrm{d}x_0$ and apply integration by parts again on the inner integral. After recombining, we get
\[
L_t(\theta) = \mathbb{E}_{x_0 \sim p_0}\,\mathbb{E}_{x_t | x_0}\left[\|s_\theta(x_t, t) - \nabla_{x_t} \log p_t(x_t | x_0)\|^2\right] + C,
\]
where $C$ is independent of $\theta$. The crucial point: the intractable marginal score $\nabla \log p$ has been replaced by the conditional score $\nabla \log p_t(\cdot | x_0)$, which is explicitly computable. Since $x_t | x_0 \sim N(\alpha_t x_0,\, \sigma_t^2 I)$, we have $\nabla_{x_t} \log p_t(x_t|x_0) = -(x_t - \alpha_t x_0)/\sigma_t^2 = -\xi_t/\sigma_t$ where $x_t = \alpha_t x_0 + \sigma_t \xi_t$ and $\xi_t \sim N(0,I)$. The final training objective is therefore
\[
\min_\theta\; \mathbb{E}_t\,\mathbb{E}_{x_0}\,\mathbb{E}_{\xi_t}\left[\left\|s_\theta(\alpha_t x_0 + \sigma_t \xi_t,\, t) - \frac{\xi_t}{\sigma_t}\right\|^2\right],
\]
which only requires samples from $p_0$ and standard Gaussians.
\end{exam}